\sgaStableTarget{sga6.front.introduction}\section*{\underline{Introduction}}
The aim of the present Seminar is to develop a \underline{global} intersection theory, and a Riemann--Roch formula, for arbitrary schemes. The reader will find a description of the Seminar's program in Exposé 0. The possibility, in principle, of proving a Riemann--Roch formula for general \underline{regular} schemes, along the lines of the well-known report of Borel--Serre in 1958, was clear from that time, at least for a projective morphism. The extension to the general case is less evident; the program of which it forms a part (and which is partially carried out in our seminar) dates from 1960. As was also the case for the duality theory of coherent sheaves (cf.\ R.~Hartshorne, Residues and Duality, Lecture Notes in Mathematics n\textsuperscript{o}~20, Springer), an essential tool for formulating a satisfactory theory is Verdier's theory of \underline{derived categories}, knowledge of which is indispensable for understanding the Seminar.

Apart from this theory, studying the Seminar requires scarcely more than a general knowledge of the foundations of Algebraic Geometry, as set out in EGA, Chapters I, II, III; in addition, we shall have only occasional references to make to EGA IV, for certain facts concerning flat or smooth morphisms, which for the most part also appear in the first exposés of SGA~1. Finally, to develop certain results with all the generality desirable for applications, we sometimes use the notions of a \underline{ringed site} and a \underline{ringed topos}, for which we refer to SGA~4, Exposés I to IV. A reader unfamiliar with the language of sites and topoi may everywhere replace these objects by ordinary topological spaces, the objects of the topos then being replaced by open subsets of these spaces; nevertheless, we advise the reader, preferably, to assimilate the language of topoi, which provides an extremely convenient principle of unification.

The fundamental notion for the theory presented here is that of a \underline{perfect complex of Modules}, and its various variants, developed in Exposés I, II, III. It seems clear that these notions, which are also important in other contexts in Algebraic Geometry (notably for formulas of Lefschetz--Verdier type in cohomology with discrete (SGA 5) or coherent coefficients), will also have their part to play outside Algebraic Geometry, notably for the formulation of a com%
